\chapter{FPGA 보드와 이야기하기: 호스트와 전시 화면}\label{ch:31-fpga-host}
\begingroup\setlength{\emergencystretch}{3em}% long identifiers: allow looser lines in this chapter only

\ref{ch:fpga-array}장은 엔진의 레지스터 버스와 레지스터 맵, 그리고 버스를 칩 밖으로 내보내는 두 브리지를
보드 쪽에서 보았다. 이 장은 반대편, PC(또는 PYNQ 보드의 ARM)에서 도는 Python 프로그램 \code{rtl/fpga/host}의
\code{bombe\_fpga} 패키지를 읽는다. 하는 일은 셋이다. 엔진에 구동을 맡기고 정지를 거둔다. 그 정지를 Rust 코어의
답과 대조한다. 그리고 전시에서 관람객이 보는 화면을 그린다.

\begin{center}
\small
\begin{tabular}{lll}
\toprule
파일 & 하는 일 & 이 장 \\
\midrule
\code{transport.py} & 32비트 읽기와 쓰기 하나씩: UART, MMIO & \ref{sec:h-transport}절 \\
\code{driver.py} & 레지스터 단어 만들기, 엔진 한 대 다루기 & \ref{sec:h-config}--\ref{sec:h-drain}절 \\
\code{farm.py} & 로터 순서를 여러 엔진에 나눠 주기 & \ref{sec:h-farm}절 \\
\code{check.py} & Rust 코어의 정지와 대조 & \ref{sec:h-check}절 \\
\code{fake.py} & 보드 없이 시험하는 가짜 보드 & \ref{sec:h-fake}절 \\
\code{visual.py} & 전시 화면(Tk) & \ref{sec:h-visual}절 \\
\code{cli.py} & 명령줄: \code{info}, \code{run}, \code{check}, \code{demo} & 곳곳에 \\
\bottomrule
\end{tabular}
\end{center}

\section{운반층: 읽기 하나, 쓰기 하나}\label{sec:h-transport}

\begin{lstlisting}[language=Python,caption={\code{transport.py}: 운반층의 약속.},label={lst:h-transport}]
class Transport:
    """Abstract 32-bit register bus (byte addresses, word aligned)."""

    def read32(self, addr: int) -> int:
        raise NotImplementedError

    def write32(self, addr: int, value: int) -> None:
        raise NotImplementedError
\end{lstlisting}

\Why 엔진을 다루는 데 필요한 것은 ``이 주소에 이 값을 써라''와 ``이 주소의 값을 읽어 달라''의 두 가지뿐이다
(\ref{sec:e-bus}절). 그 둘을 어떻게 나르는지는 보드마다 다르다. Nexys Video는 직렬선(UART)으로, PYNQ는 ARM이
메모리 주소를 읽고 쓰는 것(MMIO)으로 나른다. 위층의 코드가 이 두 함수만 쓰게 하면, 나르는 방법을 바꿔도 위층은
한 줄도 바뀌지 않는다. 시험에서는 가짜 보드(\ref{sec:h-fake}절)가 같은 두 함수를 가진다.

\General 전기 제품의 플러그와 콘센트가 이런 약속이다. 냉장고는 발전소가 무엇으로 전기를 만드는지 모른다.
콘센트의 모양과 전압만 맞으면 된다. 프로그램에서는 이것을 \term{인터페이스}{interface}라 하고, 운영체제가
파일, 네트워크, 장치를 모두 ``읽기와 쓰기''로 다루게 하는 것도 같은 생각이다.

\Ours 구현은 셋이다. \code{UartTransport}는 PC에서 Nexys Video와, \code{MmioTransport}는 PYNQ의 ARM에서
\code{pynq.MMIO}로, 그리고 가짜 보드는 Python 안에서 레지스터 맵을 흉내 낸다. \code{MmioTransport}는 네 줄짜리다.
ARM에게 엔진의 레지스터는 메모리의 한 구역이므로 \code{mmio.read(addr)}가 곧 레지스터 읽기다.

\paragraph{직렬선의 틀.} UART 쪽은 \ref{sec:e-bridge}절의 바이트 형식을 그대로 만든다.

\begin{lstlisting}[language=Python,caption={\code{transport.py}: \code{UartTransport}의 읽기와 쓰기.},label={lst:h-uart}]
def _exchange(self, frame: bytes, n_reply: int, retries: int | None = None) -> bytes:
    last = None
    for _ in range((self.retries if retries is None else retries) + 1):
        self.ser.write(frame)
        reply = self.ser.read(n_reply)
        if len(reply) == n_reply:
            return reply
        last = reply
        self.resync()
    raise TransportError(f"no/short reply ({len(last)} of {n_reply} bytes) for frame {frame.hex()}")

def read32(self, addr: int) -> int:
    # reads are idempotent except STOP_REG (offset 0x020 of an engine, which pops the FIFO):
    # that one is never retried, a lost reply is reported as an error instead of losing a stop.
    pop = addr < 0xF000 and (addr & 0xFFF) == 0x020
    reply = self._exchange(b"R" + struct.pack("<H", addr), 4, 0 if pop else None)
    return struct.unpack("<I", reply)[0]

def write32(self, addr: int, value: int) -> None:
    reply = self._exchange(b"W" + struct.pack("<HI", addr, value & 0xFFFFFFFF), 1)
    if reply != b"\x06":
        raise TransportError(f"bad write ack {reply.hex()}")
\end{lstlisting}

\code{struct.pack("<HI", addr, value)}은 주소를 2바이트, 값을 4바이트로, 작은 자리부터(\code{<}, little-endian)
늘어놓는다. 쓰기는 \code{'W'}와 이 6바이트로 7바이트, 읽기는 \code{'R'}과 주소로 3바이트다. 요청 하나를 보내고 답을
다 받을 때까지 다음 요청을 보내지 않는다(\emph{lock-step}). 브리지가 답하는 동안 들어온 바이트를 버리기 때문이다.

답이 모자라게 오면(선에서 바이트를 잃었거나 보드가 늦었으면) \code{resync}를 하고 다시 보낸다. \code{resync}는
15밀리초를 기다린 뒤 받은 바이트를 모두 버린다. 브리지는 바이트 사이가 약 10밀리초 넘게 벌어지면 받던 명령을
버리므로(\ref{sec:e-bridge}절), 15밀리초를 쉬면 브리지가 반쯤 받은 틀은 확실히 사라지고 양쪽이 처음부터 다시
시작한다.

\section{다시 묻지 말아야 할 읽기}\label{sec:h-idempotent}

위 코드에서 가장 중요한 줄은 \code{pop = ...}이다. 읽기는 원래 몇 번을 되풀이해도 결과가 같다. 그래서 답을
못 받으면 다시 물으면 된다. 그런데 \code{STOP\_REG}(엔진 안의 주소 \code{0x020})는 읽으면 정지 FIFO에서 맨 앞의
정지가 빠진다(\ref{sec:e-bus}절). 답이 오는 길에서 사라졌는데 다시 물으면, 엔진은 \emph{다음} 정지를 내준다.
첫 정지는 아무도 모르게 사라진다.

\Why 같은 요청을 여러 번 보내도 한 번 보낸 것과 결과가 같은 동작을 \term{멱등}{idempotent}하다고 한다. 통신이
불확실할 때 ``모르겠으면 다시 보낸다''는 멱등한 동작에만 안전하다. 그래서 운반층은 \code{STOP\_REG}만 재시도
횟수를 0으로 하고, 답을 잃으면 정지를 잃은 채 계속 가지 않고 오류를 낸다. 조용히 틀리는 것보다 소리 내어 멈추는
편이 낫다.

\General 은행 자동화기기를 생각하면 된다. ``잔액 조회''는 응답이 끊겨도 다시 누르면 그만이다. ``출금''은 다르다.
응답이 끊겼다고 다시 누르면 두 번 빠져나갈 수 있다. 그래서 은행의 출금 요청에는 거래 번호가 붙어 같은 요청을
두 번 처리하지 않게 한다.

\Ours 시험 \code{test\_stop\_reg\_read\_is\_never\_retried}가 이것을 확인한다. 바이트 단위의 가짜 직렬선
(\code{FakeSerial})이 답을 하나 버리게 해 두고 \code{STOP\_REG}를 읽으면, 요청이 정확히 한 번만 나가고
\code{TransportError}가 나야 한다. 다른 주소에서는 답을 한 번 버려도 재시도로 같은 값을 읽고, 다섯 번 버리면
오류가 나는지를 \code{test\_uart\_framing\_and\_retry}가 본다.

\begin{figure}[ht]
\centering
\begin{tikzpicture}[x=1cm,y=-1cm]
  \node[blkg,minimum width=22mm] (pc) at (0,0) {PC};
  \node[blkn,minimum width=22mm] (br) at (8,0) {UART 브리지};
  \draw[cgray] (0,0.4) -- (0,7.4); \draw[cgray] (8,0.4) -- (8,7.4);
  \draw[->] (0,1.0) -- node[above,font=\scriptsize\ttfamily] {57 | 08 10 | DD CC BB AA} (8,1.3);
  \draw[<-] (0,1.9) -- node[above,font=\scriptsize\ttfamily] {06} (8,1.6);
  \node[note,anchor=east] at (-0.2,1.4) {쓰기 7바이트};
  \draw[->] (0,2.6) -- node[above,font=\scriptsize\ttfamily] {52 | 08 10} (8,2.9);
  \draw[<-] (0,3.5) -- node[above,font=\scriptsize\ttfamily] {DD CC BB AA} (8,3.2);
  \node[note,anchor=east] at (-0.2,3.0) {읽기 3바이트};
  \draw[->] (0,4.2) -- node[above,font=\scriptsize\ttfamily] {52 | 08 10} (8,4.5);
  \draw[<-,dashed,cred] (4.6,4.95) -- (8,4.8);
  \node[cred,font=\small] at (4.4,4.95) {$\times$};
  \node[note,anchor=east,text=cred] at (-0.2,4.6) {답을 잃음};
  \draw[decorate,decoration={brace,amplitude=4pt,mirror}] (-0.15,5.2) -- (-0.15,5.9) node[midway,left=5pt,note] {15 ms 쉼};
  \draw[->] (0,6.2) -- node[above,font=\scriptsize\ttfamily] {52 | 08 10} (8,6.5);
  \draw[<-] (0,7.1) -- node[above,font=\scriptsize\ttfamily] {DD CC BB AA} (8,6.8);
  \node[note,anchor=west,align=left] at (8.3,5.6) {재시도\\(\code{STOP\_REG}만은\\재시도하지 않고 오류)};
\end{tikzpicture}
\caption{직렬선 위의 쓰기, 읽기, 재시도. 엔진 1의 \code{ORDER}(주소 \code{0x1008})에 \code{0xAABBCCDD}를
쓰고 읽는 예로, 바이트는 시험 \code{test\_uart\_framing\_and\_retry}의 것이다. 주소와 값은 작은 자리부터 나간다.}
\label{fig:g-h-uart}
\end{figure}

\section{구동 하나를 꾸린다}\label{sec:h-config}

엔진 한 대가 로터 순서 하나를 돌게 하려면 \code{ORDER}, \code{OPTIONS}, \code{EDGE[i]}, \code{EDGE\_EN}을
써야 한다. 그 값은 32비트 단어 안의 비트 자리에 수를 채운 것이다.

\begin{lstlisting}[language=Python,caption={\code{driver.py}: 레지스터 단어 만들기.},label={lst:h-encode}]
def encode_order(order) -> int:
    left, middle, right = (ROTOR_IDS[r] for r in order.rotors)
    word = REFLECTOR_IDS[order.reflector] | left << 8 | middle << 12 | right << 16
    if order.greek:
        word |= 1 << 2 | ROTOR_IDS[order.greek] << 4
    return word

def config_from(menu, order, diagonal_board=True, turnover_at=None, max_edges=32) -> EngineConfig:
    if len(menu.edges) > max_edges:
        raise EdgeLimitError(
            f"menu has {len(menu.edges)} edges but the FPGA supports {max_edges}: use menu.select({max_edges})")
    test, wire = ord(menu.test_letter) - 65, ord(menu.input_wire) - 65
    opt = test | wire << 8 | int(diagonal_board) << 16
    if turnover_at is not None:
        opt |= 1 << 17 | (turnover_at & 0xFF) << 24
    edges = tuple((ord(e.plain) - 65) | (ord(e.cipher) - 65) << 8 | e.pos << 16 for e in menu.edges)
    return EngineConfig(encode_order(order), opt, edges)
\end{lstlisting}

\code{x << 8}은 $x$를 여덟 자리 위로 올리는 것, \code{|}는 비트를 겹치는 것이다. 각 값이 자기 자리에만 비트를
가지므로 겹쳐도 서로 섞이지 않는다(그림~\ref{fig:g-h-words}). 받는 회로는 반대로 자리를 잘라 읽는다
(\code{order\_q[11:8]}이 왼쪽 로터).

\begin{figure}[ht]
\centering
\begin{tikzpicture}[x=3.6mm,y=1cm]
  \def\fld#1#2#3#4#5{% y, hi, lo, label, color
    \fill[#5] ({31-#2},#1) rectangle ({32-#3},#1+0.62);
    \draw ({31-#2},#1) rectangle ({32-#3},#1+0.62);
    \node[font=\tiny] at ({(63-#2-#3)/2},#1+0.31) {#4};}
  \foreach \yy/\nm in {0/EDGE[i], 1/OPTIONS, 2/ORDER} {
    \draw[cgray!60] (0,\yy) rectangle (32,\yy+0.62);
    \node[anchor=east,font=\scriptsize\ttfamily] at (-0.3,\yy+0.31) {\nm};
  }
  \fld{2}{1}{0}{반사}{cbluebg}  \fld{2}{2}{2}{G}{cyellowbg} \fld{2}{7}{4}{그리스}{cyellowbg}
  \fld{2}{11}{8}{왼쪽}{cgreenbg} \fld{2}{15}{12}{가운데}{cgreenbg} \fld{2}{19}{16}{오른쪽}{cgreenbg}
  \fld{1}{4}{0}{테스트}{corangebg} \fld{1}{12}{8}{입력 선}{corangebg} \fld{1}{16}{16}{D}{cbluebg} \fld{1}{17}{17}{T}{cbluebg}
  \fld{1}{31}{24}{넘어감 위치}{cyellowbg}
  \fld{0}{4}{0}{평문}{cgreenbg} \fld{0}{12}{8}{암호문}{cgreenbg} \fld{0}{23}{16}{크립 위치}{cyellowbg}
  \foreach \b in {31,24,16,8,0} { \node[font=\tiny,text=cgray] at ({31.5-\b},2.85) {\b}; }
\end{tikzpicture}
\caption{엔진에 쓰는 세 단어의 비트 자리(왼쪽이 31번 비트). G는 그리스 로터 있음, D는 대각선 보드, T는 넘어감
가설 켬이다. \code{rtl/fpga/README.md}의 레지스터 맵을 그린 것이다.}
\label{fig:g-h-words}
\end{figure}

\Why 메뉴의 선이 엔진의 한도(\code{MAX\_EDGES})보다 많으면 \code{config\_from}은 쓰기 전에 거절하고, 오류 글에
고칠 방법(\code{menu.select(...)})까지 적는다. 엔진에 넘치는 선을 쓰면 주소가 다른 레지스터를 덮거나 무시되어,
틀린 메뉴로 조용히 돈다. 하드웨어의 한계는 하드웨어에 닿기 전에, 사람이 읽을 수 있는 말로 알리는 것이 좋다.

\Ours \code{Engine.configure}는 이 단어들을 \code{ORDER}, \code{OPTIONS}, \code{EDGE[0]}, \ldots,
\code{EDGE\_EN} 순으로 쓰고 \code{start}는 \code{CTRL}에 1을 쓴다. 엔진은 시작 신호를 받을 때 설정을 읽으므로
쓰는 순서 자체는 상관없지만, 다 쓴 뒤에 시작한다는 것은 지켜야 한다. 시험 \code{test\_config\_words}가 각 자리의
값을 확인하고, \code{test\_order\_encoding\_roundtrip}은 단어를 다시 로터 순서로 풀어(\code{decode\_order}) 처음과
같은지 본다. 만드는 쪽과 푸는 쪽을 맞물려 시험하면 한쪽만 틀린 경우를 잡는다.

\section{정지를 거둔다}\label{sec:h-drain}

\begin{lstlisting}[language=Python,caption={\code{driver.py}: FIFO를 비우고, 끝났는지 본다.},label={lst:h-poll}]
def pop_stops(self) -> list:
    out = []
    while True:
        pos_word = self.rd(REG["STOP_POS"])
        if not pos_word >> 31:
            return out
        out.append(decode_stop(pos_word, self.rd(REG["STOP_REG"])))   # STOP_REG read pops

@staticmethod
def finished(status: int) -> bool:
    return bool(status & ST_DONE) and not status & ST_BUSY and not status & ST_STOP

def poll(self):
    st = self.status()
    if st & ST_OVERFLOW:
        raise RuntimeError(f"engine {self.index}: stop FIFO overflow (stops lost)")
    stops = []
    if st & ST_STOP:
        stops = self.pop_stops()
        st = self.status()           # re-read: done may have arrived while we drained
    return stops, self.finished(st)
\end{lstlisting}

정지 하나는 \code{STOP\_POS}를 읽고(31번 비트가 ``있음'') \code{STOP\_REG}를 읽어 꺼낸다. FIFO가 빌 때까지
되풀이한다. 구동이 끝났다는 판정에는 세 조건이 모두 필요하다. \code{DONE}이 켜졌고, \code{BUSY}가 꺼졌고,
\code{STOP}(FIFO에 정지가 있음)이 꺼졌다. 엔진이 마지막 위치를 끝냈어도 FIFO에 정지가 남아 있으면 아직 끝난 것이
아니다.

\Why \code{poll}이 정지를 비운 뒤 상태를 \emph{다시} 읽는 이유는 시간이 흐르기 때문이다. 처음 읽은 상태는
``아직 도는 중''이었는데, 정지를 비우는 동안 엔진이 마지막 위치를 마칠 수 있다. 다시 읽으면 그 차례에 끝난 것을
알아차린다. 그리고 끝났다는 판정에 ``FIFO가 비었음''이 들어 있으므로, 비우는 사이에 정지가 하나 더 들어왔다면 끝난
것으로 치지 않고 다음 차례에 마저 거둔다. 어느 쪽이든 정지가 남은 채로 엔진을 놓아주는 일은 없다.
\ref{sec:e-fifo}절의 역압 덕분에 엔진은 FIFO가 차면 기다리므로, 호스트가 아무리 느려도 정지는 빠지지 않는다.
호스트가 할 일은 비우는 것을 잊지 않는 것뿐이다. \code{STOP\_OVERFLOW}는 그래서 켜질 일이 없는 비트이고, 켜지면
무언가 크게 틀린 것이므로 곧바로 멈춘다.

\section{엔진 여러 대에 나눠 주기: \texorpdfstring{\code{Farm}}{Farm}}\label{sec:h-farm}

\begin{lstlisting}[language=Python,caption={\code{farm.py}: 노는 엔진에 다음 로터 순서를 준다(일부).},label={lst:h-farm}]
results = [[] for _ in orders]
pending = deque(enumerate(orders))
active = {}                      # engine index -> order index
while pending or active:
    for eng in engines:
        if eng.index not in active and pending:
            k, order = pending.popleft()
            eng.configure(config_from(menu, order, diagonal_board, turnover_at, self.fpga.max_edges))
            eng.start()
            active[eng.index] = k
    for eng in engines:
        k = active.get(eng.index)
        if k is None:
            continue
        stops, fin = eng.poll()
        results[k].extend(stops)
        if fin:
            ...
            del active[eng.index]
\end{lstlisting}

할 일(\code{pending})은 번호를 붙인 로터 순서의 줄이다. 한 바퀴마다 두 가지를 한다. 노는 엔진이 있으면 줄의 맨
앞 순서를 꾸려 시작시키고, 돌고 있는 엔진마다 정지를 거두고 끝났으면 놓아준다. 정지는 \emph{엔진}이 아니라
\emph{로터 순서의 번호} \code{k}의 칸(\code{results[k]})에 쌓는다.

\Why 번호의 칸에 쌓는 것은 \ref{sec:s-gather}절과 같은 이유다. 엔진 두 대가 어느 쪽이 먼저 끝날지는 정지 수에
따라 달라진다. 결과를 끝난 순서로 쌓으면 같은 입력에 대해 출력의 순서가 그때그때 달라지고, Rust 코어와 순서까지
비교할 수 없다. 번호의 칸에 쌓으면 엔진이 몇 대든, 어느 엔진이 어느 순서를 맡았든 결과는 \code{wheel\_orders}의
순서 그대로다.

\General 은행 창구의 한 줄 서기와 같다. 손님은 한 줄로 서고, 비는 창구가 맨 앞 손님을 부른다. 창구마다 따로 줄을
서면 느린 창구의 줄만 길어진다. 한 줄 서기는 일이 고르지 않을 때(어떤 로터 순서는 정지가 많아 거두는 데 오래
걸린다) 특히 낫다.

\Ours \ref{sec:e-array}절의 끝에서 이 프로그램을 ``막사의 감독''이라 불렀다. 감독이 하는 일이 이 몇 줄이다.
\code{FpgaBombe.run}은 \code{Farm}을 거쳐 \code{christopher.Bombe.run}과 \emph{같은 모양의 결과}, 곧 \code{Stop}의
목록을 돌려준다. 정지를 \code{Bombe.\_stop}으로 만들므로 FPGA가 낸 정지에도 \code{stop.check()}를 그대로 부를 수
있다. 하드웨어가 봄베를 돌리고, 체킹은 Rust 코어가 한다.

\section{Rust 코어와 대조한다}\label{sec:h-check}

\begin{lstlisting}[language=Python,caption={\code{check.py}: 정지 목록의 비교.},label={lst:h-compare}]
def _key(stop):
    return (str(stop.wheel_order), stop.position, stop.test_register)

def compare(fpga_stops, sw_stops) -> CheckReport:
    f, s = [_key(x) for x in fpga_stops], [_key(x) for x in sw_stops]
    return CheckReport(
        fpga_stops, sw_stops,
        only_fpga=[k for k in f if k not in set(s)],
        only_software=[k for k in s if k not in set(f)],
        same_order=f == s,
    )
\end{lstlisting}

정지 하나를 (로터 순서, 드럼 위치, 테스트 레지스터)로 나타내고, 두 목록을 세 가지로 비교한다. FPGA에만 있는
정지, 소프트웨어에만 있는 정지, 그리고 순서까지 같은지. 결과가 같지 않으면 다른 정지를 스무 개까지 적어 준다.
``다르다''만 알려 주는 것보다 \emph{무엇이} 다른지가 고장을 찾는 데 훨씬 쓸모 있다. 정지가 같은데 순서만
다르면 그것도 따로 알린다. 순서만 다르다는 것은 드럼을 넘기는 순서나 FIFO에 문제가 있다는 단서다.

\code{check}는 같은 메뉴와 로터 순서로 \code{Farm(fpga).run}과 \code{Bombe(...).run}을 돌려 \code{compare}에
넘긴다. 명령줄의 \code{python -m bombe\_fpga -{}-uart COM5 check -{}-ciphertext ... -{}-crib ... -{}-offset ...}이
이것을 부르고, 다르면 종료 코드 1로 끝난다. 다른 프로그램이 이 명령을 불러 합격과 불합격을 판정할 수 있다.

\Ours 보드 없이 이 경로 전체를 시험한다. \code{test\_farm\_matches\_software}는 가짜 보드(\ref{sec:h-fake}절)의 엔진
수와 FIFO 깊이를 (1대, 1칸), (2대, 4칸), (4대, 64칸)으로 바꿔 가며 로터 순서 다섯 개를 돌리고, 결과가 소프트웨어와
같은지 본다. FIFO가 한 칸이면 엔진이 거의 매 정지마다 기다려야 하므로 역압의 경로가 시험된다.
\code{test\_full\_run\_over\_uart\_fake}는 그 위에 바이트 단위의 가짜 직렬선까지 얹어, 바이트를 주고받는 것만으로
같은 결과가 나오는지 본다. 이 시험들은 호스트 프로그램의 논리를 확인한다. 실제 보드에서의 대조는 보드를 확보한
뒤의 일이다(\ref{sec:e-results}절).

\section{전시 모드: 사람의 속도로}\label{sec:h-exhibit}

엔진은 위치 하나를 15클럭 안팎, 50MHz에서 0.3마이크로초에 끝낸다. 로터 순서 하나가 5밀리초다. 사람의 눈에는
아무 일도 일어나지 않은 것과 같다. FPGA에는 돌아가는 드럼도 딸깍거리는 릴레이도 없으므로, 전시에서는 엔진이
하는 일을 \emph{화면으로} 보여 줘야 하고, 그러려면 엔진을 사람이 따라갈 수 있는 속도로 늦추고 안을 들여다볼 창을
내야 한다. 엔진의 둘째 판(\code{PARAMS}의 판 번호 2)에 더한 레지스터가 그것이다.

\begin{center}
\small
\begin{tabularx}{\linewidth}{llX}
\toprule
주소 & 이름 & 내용 \\
\midrule
\code{0x000} 비트 2 & \code{CTRL.STEP} & 붙잡힌 엔진을 한 위치만 보낸다 \\
\code{0x004} 비트 4 & \code{STATUS.HOLD} & \code{STEP}을 기다리는 중 \\
\code{0x028} & \code{PACE} & 위치마다 더 쉬는 클럭 수(0이면 최고 속도, 기본값) \\
\code{0x02C} & \code{DEMO} & 비트 0 한 걸음씩(위치마다 붙잡음), 비트 1 정지에서 붙잡음 \\
\code{0x030} & \code{DRUM} & 지금의 드럼 위치 \\
\code{0x034} & \code{TESTREG} & 지금의 테스트 레지스터 \\
\code{0x038}, \code{0x03C} & \code{SNAP\_POS}, \code{SNAP\_ITERS} & 스냅숏의 위치(있음, 정지였음 비트 포함)와 걸린 걸음 수 \\
\code{0x200} $+ 4c$ & \code{SNAP[c]} & 케이블 $c$의 26비트 스냅숏($c = 0,\ldots,25$) \\
\bottomrule
\end{tabularx}
\end{center}

\paragraph{쉬는 상태 하나.} 늦추는 방법은 상태 기계(\ref{sec:fsm}절)에 상태 하나를 더한 것이다.

\begin{lstlisting}[language=SV,caption={\code{bombe\_engine.sv}: 위치를 마친 뒤 쉬거나 기다리는 상태.},label={lst:h-pace}]
assign hold_now = single_q || (pause_stop_q && stop_cur);
...
S_PUSH: begin
  if (!fifo_full) begin
    wr_ptr <= wr_ptr + 1'b1;
    state  <= (pace_q != 32'd0 || hold_now) ? S_PACE : S_ADVANCE;
  end
end
S_PACE: begin
  // Exhibit mode: slow the sweep to human speed or wait for a STEP write.
  if (hold_now) begin
    if (step_pulse) state <= S_ADVANCE;
  end else if (pace_cnt >= pace_q) begin
    state <= S_ADVANCE;
  end else begin
    pace_cnt <= pace_cnt + 32'd1;
  end
end
\end{lstlisting}

위치 하나를 마치면(정지면 FIFO에 넣은 뒤) 원래는 곧장 드럼을 넘기는 \code{S\_ADVANCE}로 갔다. 이제 \code{PACE}가
0이 아니거나 붙잡아야 하면 \code{S\_PACE}에 들른다. 거기서 붙잡을 때가 아니면 \code{PACE}만큼 클럭을 세고 넘어가고,
붙잡을 때면 호스트가 \code{CTRL}에 \code{STEP}을 쓸 때까지 기다린다. \code{STEP}은 \code{START}처럼 한 클럭짜리
펄스다(\ref{sec:e-bus}절).

\Why 이 설계에서 지킨 원칙은 ``기본값에서는 아무것도 바뀌지 않는다''이다. \code{PACE}와 \code{DEMO}가 0이면
\code{S\_PACE}에 들르지 않으므로 위치당 클럭 수가 예전과 똑같다. 그래서 \ref{sec:e-verify}절의 시험을 그대로
다시 돌릴 수 있고, 프로젝트 계획서의 기록으로 엔진 시험 9개가 모두 통과했으며 늘어난 LUT는 약 2,500개다. 전시를
위한 기능이 기준 장치로서의 기능을 흔들면 안 된다.

\paragraph{속도를 클럭 수로.} 초당 몇 위치로 돌릴지는 호스트가 정하고, 엔진에는 클럭 수로 알린다.

\begin{lstlisting}[language=Python,caption={\code{driver.py}: 초당 위치 수를 \code{PACE}로.},label={lst:h-pacefor}]
POSITION_CYCLES = 16        # typical natural cycles per position (12..17 measured), added to PACE

def pace_for_rate(positions_per_s: float, clock_mhz: int) -> int:
    if positions_per_s <= 0:
        raise ValueError("rate must be positive")
    return max(0, round(clock_mhz * 1e6 / positions_per_s) - POSITION_CYCLES)
\end{lstlisting}

위치 하나에 쓸 수 있는 클럭은 (초당 클럭 수) $\div$ (초당 위치 수)이고, 엔진이 계산에 쓰는 약 16클럭을 뺀
나머지를 쉬면 된다. 초당 2위치이면 $50{,}000{,}000 / 2 - 16 = 24{,}999{,}984$다(\code{rtl/fpga/README.md}의 예와
같고, 시험 \code{test\_registers\_read\_back\_and\_set\_rate}가 확인한다). 엔진이 계산에 쓰는 클럭은 위치마다
12--17로 조금씩 다르므로 속도는 ``약''이다. 화면에는 엔진의 \code{POSITIONS}를 세어 잰 실제 속도를 따로 보인다.

\section{움직이는 것을 사진 찍기: 스냅숏}\label{sec:h-snap}

화면에 676가닥을 그리려면 엔진의 \code{cab\_q}(\ref{sec:e-spread}절)를 읽어야 한다. 그런데 \code{cab\_q}는 전파
도중 클럭마다 바뀌고, 호스트는 레지스터 하나를 읽는 데 직렬선 왕복이 필요하다. 26개를 읽는 동안 값은 수없이
바뀐다. 그래서 엔진은 위치 하나를 마치는 순간 \code{cab\_q}를 따로 둔 레지스터 묶음 \code{snap\_q}에 복사한다.

\begin{lstlisting}[language=SV,caption={\code{bombe\_engine.sv}: 위치를 마칠 때의 복사(\code{S\_SPREAD} 안).},label={lst:h-snapcopy}]
if (stable || test_full) begin
  positions_q <= positions_q + 32'd1;
  // Snapshot for the exhibit display (plain registers, no logic on the way in).
  snap_q       <= cab_q;
  snap_pos_q   <= {g, l, m, r};
  snap_iters_q <= iters_q;
  snap_valid_q <= 1'b1;
  snap_stop_q  <= !test_full;
  ...
\end{lstlisting}

\Why 복사하는 쪽에 논리를 하나도 두지 않은 데 주의한다(``no logic on the way in''). \code{cab\_q}에서
\code{snap\_q}로 가는 길은 선뿐이고, 26개 가운데 하나를 고르는 선택 회로는 버스 쪽(읽는 쪽)에 있다. 엔진에서 가장
느린 경로는 전파 회로이고(\ref{sec:e-results}절), 그 경로 위에 무언가를 얹으면 50MHz를 지키지 못할 수 있다.
스냅숏은 그 경로 밖에 매달려 있다.

스냅숏의 뜻은 정지냐 아니냐에 따라 다르다. 정지인 위치에서는 전류가 끝까지 퍼진 고정점이므로 소프트웨어의
\code{Bombe.trace}와 같아야 하고, 엔진 시험 \code{snapshot\_matches\_reference\_trace}가 기준 벡터
(\ref{sec:py-vectors}절)와 676비트를 비교한다. 정지가 아닌 위치에서는 테스트 레지스터가 가득 차는 순간 엔진이 그
위치를 버리므로(조기 종료), 스냅숏은 퍼지다 만 상태다. 화면에는 그것도 보여 줄 만하다. 틀린 위치에서 전류가 순식간에
번지는 모습이기 때문이다.

\paragraph{일관되게 읽기.} 스냅숏도 다음 위치를 마치는 순간 바뀐다. 26개를 읽는 도중에 바뀌면 앞의 몇 케이블은
이전 위치, 뒤의 케이블은 다음 위치의 값이 섞인다.

\begin{lstlisting}[language=Python,caption={\code{driver.py}: 앞뒤로 \code{POSITIONS}를 읽어 확인한다.},label={lst:h-snapshot}]
def snapshot(self, retries: int = 4):
    for _ in range(retries):
        before = self.rd(REG["POSITIONS"])
        word = self.rd(REG["SNAP_POS"])
        if not word >> 31:
            return None
        cables = tuple(self.rd(REG["SNAP"], c) & 0x3FFFFFF for c in range(26))
        iters = self.rd(REG["SNAP_ITERS"]) & 0xFFFF
        if self.rd(REG["POSITIONS"]) == before:
            return Snapshot(decode_drums(word), before, iters, bool(word >> 30 & 1), cables)
    raise RuntimeError("snapshot kept changing: slow the sweep (PACE) or single-step")
\end{lstlisting}

읽기 전과 뒤에 위치 계수기 \code{POSITIONS}를 읽는다. 두 값이 같으면 그 사이에 위치를 하나도 마치지 않았으므로
스냅숏도 바뀌지 않았다. 다르면 버리고 다시 읽는다. 네 번 해도 안 되면 엔진이 너무 빠른 것이니 늦추라는 오류를 낸다.
엔진을 붙잡아 둔 동안(한 걸음씩, 정지에서 멈춤)은 스냅숏이 바뀌지 않으므로 한 번에 성공한다.

\General 자동차의 주행 거리계를 사진 없이 손으로 옮겨 적는다고 하자. 앞자리부터 적는 사이에 끝자리가 9에서 0으로
넘어가면 999를 적다가 1000이 되어 엉뚱한 수를 적는다. 흔한 해결이 ``적기 전과 적은 뒤에 한 번씩 보고, 같으면
믿는다''이다. 32비트 처리 장치가 64비트 계수기를 두 번에 나눠 읽을 때 쓰는 방법과 같다.

\section{보드 없이: 가짜 보드}\label{sec:h-fake}

\code{fake.py}는 \code{bombe\_array}를 레지스터 수준에서 흉내 내는 Python 객체를 만든다. \code{read32}와
\code{write32}를 가지므로 운반층의 자리에 그대로 들어간다(\ref{sec:h-transport}절). 흉내는 두 단계다.

\begin{itemize}
\item \code{FakeArray}는 레지스터 맵만 흉내 낸다. \code{START}가 오면 등록된 함수 \code{sweep(order, options,
edges)}를 불러 정지 목록을 받고, \code{STATUS}를 몇 번 읽어야 끝나게 해서 호스트의 기다리는 고리를 시험한다.
\item \code{FakeDemoArray}의 엔진 \code{DemoEngine}은 \emph{시계에 따라} 위치를 넘긴다. 마지막으로 본 시각부터
흐른 시간 $\times$ (\code{PACE}로 정한 속도)만큼 위치를 끝낸 것으로 치고, 그 사이의 정지를 FIFO에 넣는다.
\code{DEMO}의 두 비트와 \code{STEP}, 스냅숏 레지스터도 흉내 낸다. 시험에서는 시계를 가짜로 바꿔 끼워, 0.49초에는
위치가 1이고 0.51초에는 2인지 같은 것을 정확히 확인한다(\code{test\_pace\_in\_time\_and\_single\_step}).
\end{itemize}

\code{sweep}과 스냅숏은 Rust 코어로 채운다(\code{software\_sweep}, \code{software\_snapshot}).

\begin{lstlisting}[language=Python,caption={\code{fake.py}: 레지스터를 풀어 소프트웨어 봄베를 돌린다(일부).},label={lst:h-swsweep}]
def sweep(order_word, options_word, edges):
    diagonal = bool(options_word >> 16 & 1)
    turnover = options_word >> 24 & 0xFF if options_word >> 17 & 1 else None
    want = tuple((ord(e.plain) - 65) | (ord(e.cipher) - 65) << 8 | e.pos << 16 for e in menu.edges)
    assert edges == want, "edge registers do not encode the menu"
    assert options_word & 31 == ord(menu.test_letter) - 65
    bombe = Bombe(model, diagonal).plug_up(menu)
    return [(s._pos, s.test_register) for s in bombe.run(decode_order(order_word), turnover_at=turnover)]
\end{lstlisting}

\Why 가짜 보드는 레지스터에 \emph{쓰인 단어}를 다시 풀어서 봄베를 돌린다. 호스트가 메뉴를 넘기는 데 쓴 파이썬
객체를 슬쩍 가져다 쓰지 않는다. 그래서 호스트가 단어를 잘못 만들면(비트 자리가 하나 밀리면) 가짜 보드도 틀린
메뉴로 돌고, 대조에서 드러난다. 4--5행은 더 직접적으로, 쓰인 선 단어가 메뉴와 같은지 확인한다. 가짜가 쓸모 있으려면
진짜와 \emph{같은 입구}만 쓰게 해야 한다.

\General 비행기 조종사는 실제 비행기를 타기 전에 모의 비행 장치로 훈련한다. 모의 장치는 진짜 조종석과 같은 계기와
손잡이를 가졌고, 그 뒤의 비행기만 계산으로 대신한다. 조종석이 진짜와 다르면 모의 훈련이 소용없다.

\Ours 가짜 보드는 시험만이 아니라 전시 준비에도 쓴다. 보드가 없어도 \code{python -m bombe\_fpga demo}가 가짜 보드
위에서 전시 화면을 띄운다. 이때의 메뉴는 호스트 시험과 같은 예제 전문(로터 IV-I-V로 암호화한 기상 통보)에서 크립
\texttt{WETTERVORHERSAGE}로 만든 것이고 로터 순서 \texttt{B IV-I-V}를 돈다. \code{-{}-synthetic}을 주면 소프트웨어
봄베 없이 무작위처럼 보이는 선을 그려 화면만 시험한다. 가짜 보드의 시간은 실제 엔진의 시간이 아니다. 실보드에서의
시연은 아직 하지 않았다(\code{docs/PLAN.md}).

\section{전시 화면}\label{sec:h-visual}

\code{visual.py}는 둘로 나뉜다. \code{DemoController}는 엔진과 이야기하는 논리 전부를 담고 화면 라이브러리를
전혀 쓰지 않는다. \code{DemoApp}은 Python 표준 라이브러리의 Tk로 창을 그린다.

\Why 이렇게 나누면 논리를 창 없이 시험할 수 있다. 시험 \code{test\_controller\_\allowbreak menu\_poll\_and\_controls}와
\code{test\_controller\_\allowbreak measured\_rate}는 가짜 보드와 가짜 시계 위에서 컨트롤러의 버튼 동작과 속도 측정을 확인한다.
창을 띄우는 시험은 사람이 보고 판단해야 하지만, 논리의 시험은 컴퓨터가 판정한다. 창 쪽에는 \code{-{}-smoke FRAMES}가
있어서 지정한 장 수만큼 그리고 스스로 닫는다. 적어도 그리다가 멈추지는 않는지 자동으로 확인할 수 있다.

\begin{figure}[ht]
\centering
\begin{tikzpicture}[x=1cm,y=1cm,font=\scriptsize]
  \draw[fill=black!85,draw=cgray] (0,0) rectangle (14,7.2);
  \node[anchor=west,text=white,font=\small\bfseries] at (0.2,6.9) {FPGA BOMBE};
  \fill[black!65] (0.2,6.0) rectangle (13.8,6.6);
  \node[text=white] at (7,6.3) {배너: \texttt{searching ... KMX} / \texttt{STOP at ETD  E is steckered to V}};
  \foreach \i/\lab in {0/LEFT,1/MIDDLE,2/RIGHT} {
    \draw[cgray,fill=black!70] (0.2+\i*1.05,3.5) rectangle (1.1+\i*1.05,5.7);
    \node[text=white!70] at (0.65+\i*1.05,5.5) {\lab};
    \draw[cred,line width=0.8pt] (0.25+\i*1.05,4.4) rectangle (1.05+\i*1.05,4.8);
  }
  \node[text=white!70,anchor=west] at (0.2,3.25) {드럼 글자 띠};
  \foreach \ang/\lt in {90/E,30/I,-30/W,-90/A,-150/L,150/T} { \node[circle,draw=white!70,fill=black!70,text=white,inner sep=1pt,minimum size=4mm] (n\lt) at ($(1.75,1.55)+(\ang:1.15)$) {\lt}; }
  \draw[white!50] (nE)--(nI)--(nW)--(nA)--(nE); \draw[white!50] (nE)--(nL)--(nT);
  \node[circle,draw=cred,fill=cred,inner sep=1pt,minimum size=4mm,text=white] at (nE) {E};
  \node[text=white!70,anchor=west] at (0.2,0.2) {메뉴 그래프};
  \foreach \c in {0,...,12} { \foreach \w in {0,...,12} {
      \pgfmathsetmacro{\cellon}{mod(\c*7+\w*3,5)<2 ? 1 : 0}
      \ifdim\cellon pt>0.5pt \fill[corange] (4.0+\w*0.36,5.55-\c*0.36) rectangle +(0.3,0.3);
      \else \fill[black!65] (4.0+\w*0.36,5.55-\c*0.36) rectangle +(0.3,0.3); \fi } }
  \node[text=white!70] at (6.35,0.65) {$26\times26$ 선 격자(일부), 잔상};
  \node[text=white!70,anchor=west] at (9.2,5.65) {테스트 레지스터 램프 26};
  \foreach \i in {0,...,20} { \fill[corange] ({9.35+mod(\i,7)*0.42},{5.2-floor(\i/7)*0.42}) circle (0.14); }
  \foreach \i in {21,...,25} { \fill[black!55] ({9.35+mod(\i,7)*0.42},{5.2-floor(\i/7)*0.42}) circle (0.14); }
  \node[text=white!70,anchor=west] at (9.2,3.5) {초당 위치 수, 진행 막대};
  \fill[black!65] (9.3,3.0) rectangle (12.6,3.15); \fill[corange] (9.3,3.0) rectangle (10.7,3.15);
  \node[text=white!70,anchor=west] at (9.2,2.5) {찾은 정지 기록};
\end{tikzpicture}

\vspace{1mm}
\begin{tikzpicture}[x=1cm,y=1cm,font=\scriptsize]
  \foreach \x/\lab in {0/Start / Run,2.1/Stop,4.2/Step} { \node[draw,fill=clight,minimum width=18mm] at (\x,0) {\lab}; }
  \node[draw,fill=clight] at (6.5,0) {$\square$ pause on stop};
  \node at (8.6,0) {slow}; \draw[cgray,line width=2pt] (9.0,0) -- (12.3,0); \fill[cblue] (9.7,0) circle (0.12); \node at (12.7,0) {fast};
\end{tikzpicture}
\caption{전시 화면의 배치를 줄여 그린 것(\code{visual.py}의 \code{\_build\_static}과 \code{\_build\_controls}).
격자의 무늬는 그림을 위한 것이다.}
\label{fig:g-h-screen}
\end{figure}

\paragraph{화면의 요소.} 드럼마다 글자 띠가 굴러가듯 움직이고(가운데 빨간 틀이 지금 글자), 메뉴는 글자를 원 위에
놓은 그래프로 그리며 테스트 글자를 빨갛게 칠한다. 가운데의 $26\times26$ 격자는 스냅숏이고 가로줄이 케이블, 세로줄이
선이다. 켜진 선은 주황, 꺼진 선은 그릴 때마다 그 전 밝기의 0.75배로 어두워지므로 빨리 돌 때에는 잔상이 남아 전류가
번지는 모양이 보인다. 오른쪽에는 테스트 레지스터의 램프 26개, 잰 속도, 진행 막대, 찾은 정지의 기록이 있다. 정지를
만나면 배너가 붉어지고 ``\texttt{STOP at ETD   E is steckered to V}''처럼 드럼 글자와 스테커 후보를 적는다. 후보를 읽는
규칙은 \code{Stop.stecker\_candidates}와 같다(\ref{sec:py-classes}절).

\paragraph{속도 손잡이.} 미끄럼 손잡이의 눈금 $v$($0$--$99$)를 초당 $0.5\times10^{0.06v}$위치로 바꾼다. 눈금이
고르게 움직일 때 속도는 곱으로 늘어난다. 0이면 2초에 한 위치, 99이면 초당 약 44만 위치이고, 마지막 눈금 100은
\code{PACE} 0, 곧 최고 속도다. 사람이 느끼는 빠르기는 차이가 아니라 비율이어서, 초당 1위치와 2위치의 차이는 크게
느껴지지만 10만과 10만 1위치의 차이는 느껴지지 않는다. 소리의 크기를 데시벨로 나타내는 것과 같은 이유로 로그
눈금을 쓴다.

\paragraph{두 개의 흐름.} 창을 그리는 일과 엔진과 이야기하는 일은 서로 다른 스레드가 맡는다.

\begin{lstlisting}[language=Python,caption={\code{visual.py}: 일꾼 스레드(일부).},label={lst:h-worker}]
def _worker(self):
    while not self.stop_flag:
        try:
            while True:
                self.cmds.get_nowait()()
        except queue.Empty:
            pass
        ...
        f = self.ctrl.poll()
        with self.lock:
            f.new_stops = self.latest.new_stops + f.new_stops
            self.latest = f
        ...
        time.sleep(0.02)
\end{lstlisting}

일꾼 스레드는 20밀리초마다 엔진을 한 번 훑어(상태, 위치 수, 스냅숏, 정지, 드럼) 결과를 \code{latest}에 놓는다.
창의 스레드는 33밀리초마다(초당 약 30장) \code{latest}를 가져다 그린다. 단추를 누르면 창은 할 일을 함수로 만들어
줄(\code{cmds})에 넣기만 하고, 그 함수는 일꾼이 다음 차례에 실행한다.

\Why Tk의 창은 한 스레드에서만 다뤄야 하고, 그 스레드가 멈추면 창 전체가 얼어붙는다. 직렬선으로 스냅숏 하나를
읽는 데는 레지스터 서른 개 가까이의 왕복이 필요하다. 이 일을 창의 스레드에서 하면 그동안 창이 멈춘다. 느린 일은
일꾼에게 맡기고, 둘 사이에는 자물쇠(\code{lock})로 지킨 ``가장 최근 결과'' 하나와 명령의 줄만 둔다. 엔진에 말을
거는 것은 일꾼 하나뿐이므로 직렬선의 lock-step(\ref{sec:h-transport}절)도 저절로 지켜진다.

\General 식당의 홀 직원과 주방이 이렇게 일한다. 홀 직원은 주문서를 주방 창구에 꽂고 곧바로 손님에게 돌아간다.
요리가 끝나기를 서서 기다리지 않는다. 주방은 순서대로 주문을 처리하고 완성된 접시를 창구에 올린다.

\begin{keybox}[title={이 장의 정리}]
\begin{itemize}
\item 운반층은 \code{read32}와 \code{write32} 둘뿐이다. UART, MMIO, 가짜 보드가 같은 자리에 들어간다.
\item 읽으면 바뀌는 레지스터는 재시도하지 않는다. 멱등하지 않은 요청을 다시 보내면 정지를 잃는다.
\item 하드웨어의 한계(선의 수)는 하드웨어에 닿기 전에 사람의 말로 알린다.
\item 여러 엔진의 결과는 로터 순서의 번호 칸에 모은다. 그래야 Rust 코어와 순서까지 대조할 수 있다.
\item 전시 기능은 기본값에서 아무것도 바꾸지 않게 넣는다. 스냅숏은 가장 느린 경로 밖에 매단다.
\item 움직이는 값은 앞뒤로 계수기를 읽어 일관되게 읽는다.
\item 가짜 보드는 진짜와 같은 입구(레지스터의 단어)만 쓴다. 화면의 논리는 창과 떼어 시험한다.
\end{itemize}
\end{keybox}
\endgroup
